%%%%%%%%%%%%%%%%%%%%%% file typeinst.tex %%%%%%%%%%%%%%%%%%%%%%%%%
%
% This is the LaTeX source for the instructions to authors using
% the LaTeX document class 'llncs.cls' for contributions to
% the Lecture Notes in Computer Sciences series.
% http://www.springer.com/lncs       Springer Heidelberg 2006/05/04
%
% It may be used as a template for your own input - copy it
% to a new file with a new name and use it as the basis
% for your article.
%
% NB: the document class 'llncs' has its own and detailed documentation, see
% ftp://ftp.springer.de/data/pubftp/pub/tex/latex/llncs/latex2e/llncsdoc.pdf
%
%%%%%%%%%%%%%%%%%%%%%%%%%%%%%%%%%%%%%%%%%%%%%%%%%%%%%%%%%%%%%%%%%%%
\documentclass[runningheads,a4paper]{llncs}

%%%%%%%%%%%%%%%%%%%%%%%%%%%%

\usepackage[nodisplayskipstretch]{setspace}

%\setlength{\floatsep}{5pt}
\setlength{\textfloatsep}{14pt}

%\usepackage[
%  subtle
%  %moderate
%]{savetrees}
%------

%\usepackage{times}

%\linespread{0.97}



%\addtolength{\textfloatsep}{-0.4cm}

%\addtolength{\intextsep}{-0.3cm}

%----

%\setlength{\textfloatsep}{5pt} % Space between text and top/bottom floats
%\setlength{\floatsep}{5pt}     % Space between two consecutive floats

%\setlength{\intextsep}{5pt}    % Space above and below an in-text float

%\setlength{\textfloatsep}{7pt}

%\addtolength{\floatsep}{-0.5em}

% ROMOVE THIS IF YOU HAVE EXTRA SPACE.
\usepackage[font=small,skip=3pt]{caption}

%\captionsetup[figure]{skip=-2pt}

%\usepackage{titling}
%\setlength{\droptitle}{-3cm}

%\usepackage{dutchcal}

%\usepackage{dutchcal}

%\usepackage{pzccal}

%\usepackage[cal=boondox]{mathalfa}

\DeclareMathAlphabet{\mathpzc}{OT1}{pzc}{m}{it}

\usepackage[shortlabels]{enumitem}

\usepackage{subfiles}

\usepackage{tabularx}

\usepackage{adjustbox}

\usepackage{asymptote}

\usepackage{rotating}
%\usepackage{caption}
\usepackage[T1]{fontenc}
\usepackage{amssymb}
\usepackage{amsmath}
\usepackage{bm}
\setcounter{tocdepth}{3}
\usepackage{graphicx}
\usepackage{bussproofs}
\usepackage{supertabular}
\usepackage{longtable}
\usepackage{url}
\usepackage{afterpage}
\usepackage{tikz}
\usetikzlibrary{arrows}
\usetikzlibrary{decorations.pathmorphing}
\afterpage{\clearpage}
%\newtheorem{definition}{Definition}
%\newtheorem{proof}{Proof}
%\newproof{proof}{Proof}
%\newtheorem{mydef}{Definition}
%\newtheorem{myexample}{Example}
%\newtheorem{mytheorem}{Theorem}
%\newtheorem{mycorollary}{Corollary}
%\newtheorem{myobservation}{Observation}
%\newtheorem{mylemma}{Lemma}
\newtheorem{conj}{Conjecture}
%\newtheorem{inv}{Invariant}

\newtheorem{obs}{Observation}
%\newtheorem{proofoutline}{Proof outline}
%\newtheorem{myproof}{Proof of Proposition}
%\newtheorem{myproposition}{Proposition}
%\newtheorem{mylemma}{Lemma}
%\newtheorem{lemmaproof}{Proof of Lemma}
%\newtheorem{conjecture}{Conjecture}

%proof\_latex\_notheorem = [``proof'']

%\usepackage{amsthm}

%\usepackage{ntheorem}

\usepackage[lined,ruled,vlined,commentsnumbered]{algorithm2e}

%\usepackage{enumitem}

\usepackage{pbox}

\usepackage{stmaryrd}

\usepackage{varwidth}
\usepackage{booktabs}

\usepackage[mathscr]{euscript}

\usepackage{cancel}

%%%%%%%%%%

\begin{document}

\mainmatter  % start of an individual contribution

% first the title is needed
\title{Efficient Communication and Coordination in Multiagent Reinforcement Learning\\(Invited Talk)}

% a short form should be given in case it is too long for the running head
%\titlerunning{Efficient Coordination in MARL (Invited Talk)}

\author{Mehdi Dastani}
%\author{Neda Saeedloei}
%\affil{Towson University}
%
%\authorrunning{Neda Saeedloei and Feliks Klu\'{z}niak}
% abbreviated author list (for running head)
%
%%%% list of authors for the TOC (use if author list has to be modified)
%\tocauthor{Neda Saeedloei, Feliks Klu\'{z}niak}
%
%\institute{Towson University\\
%\email{nsaeedloei@towson.edu}

%\author{Neda Saeedloei\inst{1} \and Feliks Klu\'{z}niak\inst{2}}
%\author{Neda Saeedloei}
%
%\authorrunning{Neda Saeedloei and Feliks  Klu\'{z}niak}
%\authorrunning{Neda Saeedloei}
% abbreviated author list (for running head)
%
%%%% list of authors for the TOC (use if author list has to be modified)
%\tocauthor{Neda Saeedloei, Feliks Klu\'{z}niak}
%\tocauthor{Neda Saeedloei}
%
\institute{Utrecht University\\
  \email{M.M.Dastani@uu.nl}
%  ,
%\and
%RelationalAI\\
%\email{feliks.kluzniak@relational.ai}
}

\toctitle{Efficient Communication and Coordination in Multiagent Reinforcement Learning (Invited Talk)}
%\tocauthor{Authors' Instructions}
\maketitle

% Without this the first footnote will have number 3, since there are footnotes
% in the title.
%\setcounter{footnote}{0}

%\vspace{-1.6em}
\begin{abstract}
Multiagent reinforcement learning (MARL) is a powerful framework for learning effective policies in complex multiagent environments. A key challenge in MARL is to learn coordinated policies efficiently in a decentralised manner. Prior work has explored a range of communication and coordination mechanisms to address these challenges. In this talk, I will present our recent contributions towards improving the efficiency and performance of MARL systems through principled communication and coordination mechanisms. First, I introduce a decentralised communication scheduling approach that leverages supervised learning to construct a message estimation model. This model enables individual agents to selectively decide when sharing local information is beneficial, thereby reducing unnecessary communication. Empirical results show that this approach significantly decreases communication overhead while improving overall learning performance. Second, I present two logic-based methods for synthesising multiagent reward machines and action-masking artefacts. These methods provide formal guarantees that the resulting policies are both coordinated and safe, while also improving sample efficiency. 
\end{abstract}

\end{document}
